În prezent, aplicațiile web au devenit extrem de răspândite și sunt utilizate zilnic: fie în activitatea profesională, fie pentru comunicare, fie pentru divertisment. Această creștere semnificativă a utilizării aplicațiilor web se datorează faptului că o aplicație accesibilă din browser poate ajunge foarte aproape de nivelul și funcționalitatea programelor desktop.

În spatele aplicațiilor web moderne se află numeroase etape de planificare, decizii, activități de dezvoltare, precum și obstacole tehnice sau chiar financiare, pentru rezolvarea și implementarea cărora este responsabilă o echipă. Fiecare membru are un rol bine definit, de la formularea ideii și proiectarea UI-ului, până la dezvoltare și mentenanță.

Lucrarea își propune să prezinte procesele de planificare ale unei astfel de aplicații web: modul în care se ajunge de la idee la un produs final și ce procese au loc înainte ca dezvoltatorii să scrie chiar și o singură linie de cod sau ca designerii UX să traseze măcar o linie.

Pentru a demonstra acest lucru, este prezentată o echipă de dezvoltare „fictivă”, care proiectează și dezvoltă o aplicație web pentru jurnalizarea filmelor.

\textbf{Cuvinte cheie}: aplicații web, planificare, dezvoltare, design UX, jurnal de filme
